%%%%%%%%%%%%%%%%%%%%%%%%%%%%%%%%%%%%%%%%%
% Jiwon Kang -- Academic CV
% Self-contained single-file template (compile with pdflatex).
% Style adapted to a CVLAB-house academic CV layout.
% NOTE: resume.cls is no longer used by this file.
%%%%%%%%%%%%%%%%%%%%%%%%%%%%%%%%%%%%%%%%%

\documentclass[11pt,letterpaper]{article}

%==================== GEOMETRY & ENCODING ====================
\usepackage[T1]{fontenc}
\usepackage[utf8]{inputenc}
\usepackage[letterpaper,left=0.75in,right=0.75in,top=0.75in,bottom=0.75in]{geometry}

%==================== FONTS ====================
\usepackage{lmodern}              % math symbols stay Latin Modern (matches reference)
\usepackage{charter}              % body text = Bitstream Charter (reference font); loaded AFTER lmodern so it wins for text
\usepackage{microtype}
\usepackage[parfill]{parskip}

%==================== COLOR ====================
\usepackage{xcolor}
\definecolor{linkblue}{HTML}{1565C0}
\definecolor{cvgray}{gray}{0.45}
\definecolor{rulegray}{gray}{0.0}

%==================== TABLES / LISTS ====================
\usepackage{array}
\usepackage{tabularx}
\usepackage{enumitem}
\usepackage{ragged2e}
\usepackage[normalem]{ulem}

%==================== ICONS (graceful fallback if fontawesome5 missing) ====================
\IfFileExists{fontawesome5.sty}{\usepackage{fontawesome5}}{%
  \providecommand{\faGlobe}{[web]}\providecommand{\faGraduationCap}{[scholar]}%
  \providecommand{\faEnvelope}{[email]}\providecommand{\faGithub}{[github]}%
  \providecommand{\faLinkedin}{[linkedin]}\providecommand{\faHome}{[home]}}

%==================== SECTION HEADERS (Title-Case bold + rule) ====================
\usepackage{titlesec}
\titleformat{\section}
  {\large\bfseries}{}{0em}{}
  [\vspace{2pt}{\color{rulegray}\titlerule[0.8pt]}]
\titlespacing*{\section}{0pt}{8pt}{4pt}

%==================== HYPERLINKS (load last) ====================
\usepackage[hidelinks,breaklinks=true]{hyperref}
\hypersetup{
  colorlinks=true,
  urlcolor=linkblue,
  linkcolor=linkblue,
  citecolor=linkblue,
  pdftitle={Jiwon Kang -- Curriculum Vitae},
  pdfauthor={Jiwon Kang}
}

%==================== MACROS ====================
% Update this whenever you edit the CV content (do NOT use \today).
\newcommand{\lastupdate}{July 6, 2026}
\newcommand{\cvgap}{\hspace{1.5em}}

\newcommand{\cvheader}{%
  {\Huge\bfseries Jiwon Kang}\par
  \vspace{2pt}
  {\color{cvgray}\small Last update: \lastupdate}\par
  \vspace{6pt}
  {\small
    % \href{https://TODO.github.io}{\faGlobe\ Personal Website}\cvgap   % TODO: personal website
    \href{https://scholar.google.com/citations?user=A2PurdIAAAAJ}{\faGraduationCap\ Google Scholar}\cvgap
    \href{mailto:jiwon7258@kaist.ac.kr}{\faEnvelope\ jiwon7258@kaist.ac.kr}\cvgap
    \href{https://github.com/loggerJK}{\faGithub\ loggerJK}\cvgap
    \href{https://www.linkedin.com/in/jiwon-kang-b02b0911b/}{\faLinkedin\ LinkedIn}%
  }\par
  \vspace{10pt}
}

% Education / Experience entry: \cventry{Title}{Date}{2nd-line-left}{2nd-line-right}
% Pass {}{} for the 2nd line to omit it (robust empty test).
\makeatletter
% One entry line: left text, right text flush to the margin.
% \parfillskip=0pt so the trailing \hfill takes ALL the stretch (date reaches the right edge).
% Trailing \null (a zero-width box) keeps the last node a box, so \par never deletes the
% dangling \hfill when #2 is empty -- otherwise that single line would over-justify (interword
% glue stretched across the full width).
\newcommand{\cv@line}[2]{{\parfillskip=0pt\relax\noindent#1\hfill#2\null\par}}
\newcommand{\cventry}[4]{%
  \cv@line{\textbf{#1}}{#2}%
  \def\cv@second{#3}%
  \ifx\cv@second\@empty\else\cv@line{#3}{#4}\fi
}
\makeatother

\newlist{cvbullets}{itemize}{1}
\setlist[cvbullets]{label=\textbullet, leftmargin=1.4em, topsep=1pt,
                    itemsep=0pt, parsep=0pt, before=\small}

% Publication helpers
\newcommand{\me}[1]{\textbf{\textit{\uline{#1}}}}   % CV owner's name
\newcommand{\eqc}{\textsuperscript{*}}              % equal contribution
\newcommand{\corr}{\textsuperscript{\dag}}          % corresponding author
\newcommand{\publegend}{%
  {\color{cvgray}\small *: Equal Contribution,\quad \dag: Corresponding Author}\par
  \vspace{4pt}%
}
\newcommand{\pubgroup}[1]{\vspace{2pt}{\Large\bfseries #1}\par\vspace{2pt}}

% Reusable full link blob (Paper | Project Page | Code). Inline a subset if some URLs are unknown.
\newcommand{\publinks}[3]{%
  \small\color{linkblue}%
  \href{#1}{Paper}~\textbar~\href{#2}{Project Page}~\textbar~\href{#3}{Code}%
}

% \pubitem{label}{title}{links}{authors}{venue}{bullets}
% Title (bold) wraps in the flexible X column; links pin top-right on line 1.
% #6 = one or more \item lines describing the work; pass {} to omit
% (empty test mirrors \cventry, hence the \makeatletter wrapper for \pub@desc/\@empty).
\makeatletter
\newcommand{\pubitem}[6]{%
  \begingroup
  \setlength{\tabcolsep}{0pt}%
  \noindent
  \begin{tabularx}{\linewidth}{@{}>{\RaggedRight\arraybackslash}X @{\hspace{1.5em}} r@{}}
    \textbf{#1\ #2} & #3
  \end{tabularx}\par
  \vspace{-4pt}%
  #4\par
  \vspace{-4pt}%
  \def\pub@desc{#6}%
  \ifx\pub@desc\@empty\else
    \begin{cvbullets}#6\end{cvbullets}%
  \fi
  \vspace{-4pt}%
  \textit{#5}\par
  \endgroup
}

% to LARGE section headers
\newcommand{\largesection}[1]{\vspace{5pt} \section{\LARGE#1} \vspace{5pt}}

\makeatother

%==================== DOCUMENT ====================
\begin{document}
\pagestyle{empty}

\cvheader

%-------------------- RESEARCH INTERESTS --------------------
\largesection{Research Interests}

% I am broadly interested in generative models, world modeling, and multimodal learning. 

My research focuses on generative models, particularly diffusion/flow-based methods and their applications across vision, 3D, and multimodal domains.
Specifically:
\begin{cvbullets}
  \vspace{-5pt}
  \item Diffusion \& Generative AI (Guidance, Sampling, Image Generation)
  \item Unified Multimodal Models
  \item 3D Vision (Novel View Synthesis, NeRF, Gaussian Splatting)
\end{cvbullets}

%-------------------- EDUCATION --------------------
\largesection{Education}
\cventry{M.S. in Artificial Intelligence}{Mar. 2025 -- Feb. 2027 (Expected)}
        {Kim Jaechul Graduate School of AI, KAIST}{Advisor: \href{https://cvlab.kaist.ac.kr}{Seungryong Kim} (CVLAB)}
\vspace{2pt}
\cventry{B.S. in Computer Science and Engineering}{Mar. 2019 -- Feb. 2025}
        {Korea University}{}
\begin{cvbullets}
  \item GPA: 4.24 / 4.5
  \item Double Major in Statistics
  \item Undergraduate Research Intern @ \href{https://cvlab.kaist.ac.kr}{CVLAB}
        (Sep. 2022 -- Feb. 2025), Advisor: Seungryong Kim
\end{cvbullets}

%-------------------- EXPERIENCE --------------------
\largesection{Experience}
\cventry{Kakao, Seoul}{Mar. 2026 -- May 2026}{Research Intern}{}
\begin{cvbullets}
  \item Worked on Text-to-Speech research, repurposing Whisper as a continuous tokenizer for autoregressive diffusion-based TTS model.
\end{cvbullets}

%-------------------- PUBLICATIONS --------------------
\largesection{Publications}
\publegend

\pubgroup{International Conference}

\pubitem{[C4]}
  {Transferability Between Understanding and Generation in Unified Multimodal Models}
  {\publinks{https://arxiv.org/abs/2607.04423}{https://cvlab-kaist.github.io/UMM_Transferability/}{https://github.com/cvlab-kaist/UMM_Transferability}}
  {\me{Jiwon Kang}\eqc, Heeji Yoon\eqc, Jaewoo Jung, Jaewon Min, Minkyeong Jeon, Biyeon Hwang, Sangwon Jung, Seungryong Kim\corr}
  {European Conference on Computer Vision (ECCV 2026).}
  {\item Investigated how understanding and generation capabilities (e.g., counting) transfer to one another within unified multimodal models. \item Showed that transferability depends on the model architecture, and proposed a practical post-training strategy for unified multimodal models that leverages transferability.}

\clearpage
\pubitem{[C3]}
  {APPLE: Attribute-Preserving Pseudo-Labeling for Diffusion-Based Face Swapping}
  {\publinks{https://arxiv.org/abs/2601.15288}{https://cvlab-kaist.github.io/APPLE/}{https://github.com/cvlab-kaist/APPLE}}
  {\me{Jiwon Kang}\eqc, Yeji Choi\eqc, JoungBin Lee, Wooseok Jang, Jinhyeok Choi, Taekeun Kang, Yongjae Park, Myungin Kim, Seungryong Kim\corr}
  {IEEE/CVF Conference on Computer Vision and Pattern Recognition (CVPR 2026).}
  {\item Prior diffusion-based face swapping methods focused on changing facial identity, failing to preserve facial attributes. \item Proposed a novel pseudo-labeling framework that preserves facial attributes while swapping identities, achieving Pareto-optimal results.}
  % {\small\color{linkblue}\href{https://arxiv.org/abs/2601.15288}{Paper}~\textbar~\href{https://cvlab-kaist.github.io/APPLE/}{Project Page}}

\pubitem{[C2]}
  {A Noise is Worth Diffusion Guidance}
  {\publinks{https://arxiv.org/abs/2412.03895}{https://cvlab-kaist.github.io/NoiseRefine/}{https://github.com/cvlab-kaist/NoiseRefine}}
  {Donghoon Ahn\eqc, \me{Jiwon Kang}\eqc, Sanghyun Lee, Jaewon Min, Wooseok Jang, Minjae Kim, Hyungwon Cho, Sayak Paul, SeonHwa Kim, Eunju Cha\corr, Kyong Hwan Jin\corr, Seungryong Kim\corr}
  {International Conference on Learning Representations (ICLR 2026).}
  {\item Showed that classifier-free guidance can be replaced by refining the initial noise. \item Enabled high-quality generation without a guidance pass, highlighting the importance of the initial noise in diffusion models.}

\pubitem{[C1]}
  {Where and How to Perturb: On the Design of Perturbation Guidance in Diffusion and Flow Models}
  {\publinks{https://arxiv.org/abs/2506.10978}{https://cvlab-kaist.github.io/HeadHunter/}{https://github.com/cvlab-kaist/HeadHunter}}
  {Donghoon Ahn\eqc, \me{Jiwon Kang}\eqc, Sanghyun Lee, Minjae Kim, Wooseok Jang, Jaewon Min, Sangwu Lee, Sayak Paul, Seungryong Kim\corr}
  {Conference on Neural Information Processing Systems (NeurIPS 2025).}
  {\item Systematically studied perturbation guidance in diffusion and flow models. \item Showed that perturbing different layers/heads in a diffusion model exhibits different effects on the generated samples, which can be systematically combined to improve user-specified attributes via the proposed framework.}

\pubgroup{Workshop}

\pubitem{[W1]}
  {Self-Evolving Neural Radiance Fields}
  {\publinks{https://arxiv.org/abs/2312.01003}{https://cvlab-kaist.github.io/SE-NeRF/}{https://github.com/cvlab-kaist/SE-NeRF}}
  {Jaewoo Jung\eqc, Jisang Han\eqc, \me{Jiwon Kang}\eqc, Seongchan Kim, Min-Seop Kwak, Seungryong Kim\corr}
  {ICCV 2025 Workshop on Wild3D: 3D Modeling, Reconstruction, and Generation in the Wild.}
  {\item Proposed a self-training framework that lets a NeRF improve itself from few input views via pseudo-label distillation.}

\pubgroup{Preprints \& Under Review}

% \pubitem{[P2]}
%   {Exploring Whisper as a Continuous Tokenizer for Text-to-Speech}
%   {} % TODO: add links
%   {\me{Jiwon Kang}, Heeji Yoon, Woo-Young Kang, Sunghun Kang, Jihwan Eom, Jaeyeol Jeon, Byungseok Roh, Seungryong Kim\corr}
%   {Under Review.}
%   {\item Repurpose Whisper encoder features as a continuous tokenizer for text-to-speech.
%    \item Work done during a research internship at Kakao.}

\pubitem{[P1]}
  {RAIN-GS: Relaxing Accurate Initialization Constraint for 3D Gaussian Splatting}
  {\publinks{https://arxiv.org/abs/2403.09413}{https://cvlab-kaist.github.io/RAIN-GS/}{https://github.com/cvlab-kaist/RAIN-GS}}
  {Jaewoo Jung\eqc, Jisang Han\eqc, Honggyu An\eqc, \me{Jiwon Kang}\eqc, Seonghoon Park\eqc, Seungryong Kim\corr}
  {Under Review.}
  {\item Relaxed the accurate point-cloud initialization requirement of 3D Gaussian Splatting, enabling training from sparse or random initialization.}

%-------------------- HONORS (TODO: confirm accuracy before un-commenting) --------------------
% \section{Honors \& Awards}
% \cventry{Outstanding Paper Award (Silver), IPIU 2024}{Feb. 2024}{}{}
% \cventry{1st Place, Seoul Asan Hospital HeLP Challenge (Stroke DWI Segmentation)}{Oct. 2021}{}{}

%-------------------- INVITED TALKS (TODO) --------------------
% \section{Invited Talks}
% \cventry{Talk Title}{Date}{Hosting Venue}{}

%-------------------- ACADEMIC SERVICES (TODO) --------------------
% \section{Academic Services}
% {\bfseries Conference Reviewer}
% \begin{cvbullets}
%   \item Conference Name (ACRONYM), Year
% \end{cvbullets}

%-------------------- TEACHING (TODO: confirm) --------------------
% \section{Teaching Experience}
% \begin{cvbullets}
%   \item Teaching Assistant, Computer Vision Learning Course, Samsung Electronics (Sep. 2023)
%   \item Teaching Assistant, Visual Question Answering, Samsung Electronics (Sep. 2023)
% \end{cvbullets}

%-------------------- EXTRACURRICULAR --------------------
\largesection{Extracurricular Activities}
\cventry{AIKU (Artificial Intelligence Society of Korea University)}{Sep. 2022 -- Aug. 2023}
        {Co-founder \& Vice President of academic research society for undergraduate students}{}
% \begin{cvbullets}
%   % \vspace{-5pt}
%   \item Academic research society for undergraduate students
% \end{cvbullets}

%-------------------- TECHNICAL SKILLS / COURSEWORK (TODO) --------------------
% \section{Technical Skills}
% \begin{cvbullets}
%   \item Programming: Python, C/C++
%   \item Frameworks: PyTorch, JAX
% \end{cvbullets}

\end{document}
